\scanpage{38}

% B1 | text | first five lines
Supposons qu'on ait pu associer \uline{canoniqt},
à la catégorie de Deuring-Brauer $C$ (en termes de
la seule classe d'équivalence de $C$), un objet
$T$ de $\unsure{\nu|\mathcal V|V}(C)$. Cela signifie en particulier que
nous avons relevé le foncteur canonique \inserted{de \unsure{Gr-cat|Gp-cat}}

% B2 | diagram | upper pair of opposite arrows
\diagram{\begin{tikzcd}[ampersand replacement=\&]
\underline{\mathrm{Equ}}(C,T)\arrow[r]\&\underline{\mathrm{Equ}}(C)\arrow[l,bend left=27,"{\deleted{\unsure{\pi}}\,c}"]
\end{tikzcd}}
\unsure{en}

% B3 | text | below curved arrow
Donc pour tte catégorie $\unsure{C_1|C'}$ équivalente à $C$,
on obtient un objet $T_{1\deleted{\unsure{C'}}}^{\unsure{c|\xi}}$ canonique dans $\nu(C')$.

% B4 | text | middle through quotient
On voit donc que pour tt $(A',C',T')$ comme
plus haut, \deleted{l'ens. des} tel que $(C',T')$ soit
équivalent à $(C,T)_{A'}$, les équivalences
en question font une catégorie-torseur sous
$\underline{\mathrm{Equ}}((C,T)_{A'})$ \inserted{$\simeq\underline{\mathrm{Equ}}(C',T')$} donc \deleted{définissent une catégorie torseur}
les classes d'isom d'objets font
un torseur sous \deleted{$\underline{\mathrm{Equ}}(C',T')/\inserted{\operatorname{Im}}\underline{\mathrm{Equ}}(C)$ par passage}
$\unsure{\mathrm{Brd}|\mathrm{Bord}}(C',T')$, donc définissent un torseur sous
$\mathrm{Brd}(C',T')/\operatorname{Im}\mathrm{Brd}(C)$. Le torseur ne dépend
plus du choix d'un $C$ dans la classe
de catégories envisagée\ldots\ Il faut donc expliciter
l'application composée

% B5 | diagram | foot, three aligned rows
\diagram{\begin{tikzcd}[ampersand replacement=\&,column sep=large,row sep=small]
\mathrm{Brd}(C)\arrow[r,"c"]\&\mathrm{Brd}(C,T)\arrow[r]\&\mathrm{Brd}(C',T')\\
\unsure{\beta|\rho}\arrow[r,mapsto]\&(\beta,u_\beta)\arrow[r,mapsto]\&(\beta',u_{\beta'})\\
{}_{o}P_{o}\arrow[u,phantom,"\unsure{\wr|\simeq}" description]\arrow[r,mapsto]\&N_{o/A}(P_o)\ni u_\beta\&N_{o'/A'}(P'_o)\ni(u_\beta)'\arrow[l]
\end{tikzcd}}
% archive sheet number: 37
